\section{Introduction}

For an LLM agent to be useful in practice, it has to call several tools in sequence for a single query. A request such as ``find a news article on this stock and summarize it'' requires calling \textsc{NewsTool} and then \textsc{ResearchHelper}; no amount of single-call benchmark accuracy translates into deployment value if this chaining is unreliable. Four directions have been explored to strengthen this multi-tool ability, and they differ along a cost-vs.-targetedness trade-off.

Fine-tuning on large tool-use corpora~\citep{patil2023gorilla,qin2024toolllm,zhang2024xlam} is strong but requires weight updates and has to be repaid each time the tool catalog changes. Prompt scaffolding~\citep{yao2023react} needs no training, but inference cost scales with prompt length and degrades quickly as the number of tools grows. KV-cache optimizations~\citep{liu2024kivi} cache tool schemas efficiently and reduce latency, but focus on ``how to save memory'' rather than ``what to pick'' --- they do not intervene at the level of task decisions. The fourth direction, which this paper belongs to, is \emph{activation steering}: low-rank interventions applied to attention at inference time, without touching weights or prompts. It is cheaper than fine-tuning and more targeted than prompt engineering. Table~\ref{tab:positioning} places the four approaches side by side and locates this paper inside that landscape.

\begin{table}[t]
\centering
\caption{Four approaches to multi-tool tool calling and where this paper sits. Compared along training cost, task-level targeting, and whether the operator carries emitted-tool state.}
\label{tab:positioning}
\small
\begin{tabularx}{\linewidth}{@{}lXccc@{}}
\toprule
Approach & Representative methods & Training-free & Multi-tool targeted & Emitted state \\
\midrule
Tool-aware fine-tuning & Gorilla, ToolLlama, xLAM & $\times$ & $\checkmark$ & $\checkmark$ \\
Prompt scaffolding & ReAct, ToT & $\checkmark$ & $\checkmark$ (prompt) & prompt-level \\
KV-cache optimization & KIVI, StreamingLLM & $\checkmark$ & $\times$ & $\times$ \\
Activation steering (stationary) & SEKA, AdaSEKA, ITI, CAA, PASTA & $\checkmark$ & $\times$ (single-concept) & $\times$ \\
\textbf{This paper} & Q-coverage, Layer-adaptive K+Q & $\boldsymbol{\checkmark}$ & $\boldsymbol{\checkmark}$ (sign-flipped) & \textbf{history-free approx.} \\
\bottomrule
\end{tabularx}
\end{table}

The third and fourth columns carry the point. Existing activation-steering methods share the benefit of zero training cost, but \emph{none of them carries information about which tools have already been emitted}. The closest prior method, SEKA~\citep{SEKA}, and its extension AdaSEKA~\citep{AdaSEKA}, apply a stationary key-side amplifier of the form $k' = k + \alpha P k$, where the fixed $\alpha$ and projector $P$ depend by definition on neither the decoding step nor the already-emitted tool set. On single-concept editing (CounterFact) this stationarity is a strength; on multi-tool the same property produces ``re-amplify the first tool at the second step.'' The paper's first observation is that this is not a tuning failure but a \emph{consequence of the operator form itself}, which Theorem~\ref{thm:no-memory} fixes in a one-line structural argument.

What makes the theory interesting is that it predicts where its footprint lives in data. Theorem~\ref{thm:no-memory} implies that under SEKA the \texttt{repeated\_first\_tool\_rate} must exceed that of \texttt{no\_steer}. If this prediction holds, SEKA's F1 loss on multi-tool is repetition baked into the design, and the resolution is not a more refined SEKA but ``go \emph{less} in the same direction.'' That observation leads directly to a sign-flipped operator that subtracts from the query rather than amplifying the key.

The paper has two central claims. On the \emph{tool-selection performance} axis (\S\ref{sec:exp:performance}), we show that Q-coverage and the layer-adaptive K+Q give consistent gains over SEKA and AdaSEKA on Qwen2.5-7B-Instruct for both MetaTool Subtask4 and $\tau^2$-bench retail under matched prompt/scorer/decoding, and that the gain is interpretable as a reduction of \texttt{repeated\_first\_tool\_rate} in a position-level decomposition. Basis ablation (feature-shuffled, random, PCA-of-K) separates this gain from basis quality. On the \emph{efficiency} axis (\S\ref{sec:exp:efficiency}), the method works purely as inference-time hooks, with zero training cost compared to fine-tuning, zero token overhead compared to prompt scaffolding, and storage overhead limited to a per-head basis $(L, H_{\mathrm{kv}}, d, r)$ ($\approx 2.6$ MB on Qwen 7B).

The rest of the paper follows these two axes. \S\ref{sec:related} maps the shared structural gap across stationary steering methods. \S\ref{sec:theory} compactly presents the setup and three core theorems (corollaries and the conditional proposition move to the appendix). \S\ref{sec:exp:performance} and \S\ref{sec:exp:efficiency} support the two central claims in turn. \S\ref{sec:discussion} discusses what the result implies for other stationary steering methods.
